\chapter{Regular Conditional Distribution}

\section{From conditional expectation to conditioning on \texorpdfstring{$X=x$}{X=x}}

We would like to define \emph{conditional probability} using our conditional expectation, so we can can answer questions like: what is the probability of $Y$ to happen given that $X=x$? In principle, thanks to the factorization lemma we can formalize this idea (Note that throughout this section, $(\Omega,\A,\mathbb{P})$ is a probability space).

\begin{definition}[Conditional expectation given $X=x$]
\label{def:condexp_x}
Let $Y\in L^1(\mathbb{P})$ and let $X:(\Omega,\A)\to(E,\E)$ be measurable.
Let $Z:=\mathbb{E}[Y\mid \sigma(X)]$. By the factorization there exists $\xi:E\to\mathbb{R}$ measurable with
$Z=\xi(X)$ a.s. We define
\[
\mathbb{E}[Y\mid X=x] := \xi(x).
\]
Equivalently, for $A\in\A$,
\[
\mathbb{P}(A\mid X=x) := \mathbb{E}[\Ind_A\mid X=x].
\]
\end{definition}

However, there is a subtlety. The random variable $Z$ is not uniquely determined \emph{pointwise}: any two versions agree only $\mathbb{P}$-a.s.
In other words, for each fixed $A$ there exists a null set $N_A$ (depending on $A$) such that
two valid choices of $\mathbb{P}(A\mid \F)$ may differ on $N_A$, but nowhere else.
This is harmless if we look at one single event (or finitely many), because we can always ignore a finite union of null sets.

\paragraph{Why this is a problem when conditioning on $X=x$?}
Now suppose we want to interpret conditioning on a random variable $X$.
A natural idea is:
\[
\mathbb{P}(Y\in B\mid X=x)\stackrel{?}{=}\text{``a version of }\mathbb{E}[\mathbf{1}_{\{Y\in B\}}\mid \sigma(X)]\text{ evaluated at }X=x\text{''}.
\]
For each fixed $B\in\E$, the random variable
\[
Z_B(\omega):=\mathbb{P}(Y\in B\mid \sigma(X))(\omega)
\]
exists and is $\sigma(X)$-measurable, hence can be written as $Z_B=\xi_B(X)$ a.s.\ for some measurable $\xi_B$.
So we might want to \emph{define} $\mathbb{P}(Y\in B\mid X=x):=\xi_B(x)$.

The issue is that the null set on which $Z_B\neq \xi_B(X)$ can depend on $B$.
Thus, for each $B$ we can find a ``good'' set $\Omega_B$ with $\mathbb{P}(\Omega_B)=1$ where the identity holds,
but there is no reason for the intersection $\bigcap_{B\in\E}\Omega_B$ to still have probability $1$.
And this is exactly what we would need in order to obtain a \emph{single} object
\[
x\mapsto \big(\text{a probability measure on }(E,\E)\big)
\]
that works for \emph{all} sets $B$ simultaneously.

\paragraph{Why ``one null set per event'' is too weak?}
Think of $\mathbb{P}(Y\in B\mid \sigma(X))$ as a family of random variables indexed by $B\in\E$.
Conditional expectation guarantees:
\begin{itemize}
\item for each \emph{individual} $B$, there is a well-defined (a.s.) value of $\mathbb{P}(Y\in B\mid \sigma(X))$,
\item but the exceptional set where the chosen representative is wrong may vary with $B$.
\end{itemize}
To build a conditional \emph{distribution}, however, we need a single $\omega$-wise probability measure
$B\mapsto \vartheta(\omega,B)$, i.e.\ we need the measure axioms (countable additivity, etc.) to hold
for each fixed $\omega$ (outside one global null set).
That forces us to control \emph{countably many} relations simultaneously (at least on a countable generator),
and then extend to all $B$ by a monotone class argument.
Plain conditional expectation does not give this uniform control automatically. But the concept of regular conditional distributions fix these problems.

\begin{definition}[Transition kernel / Markov kernel]
\label{def:kernel}
Let $(\Omega_1,\A_1)$ and $(\Omega_2,\A_2)$ be measurable spaces.
A map $\vartheta:\Omega_1\times \A_2\to [0,\infty]$ is a \emph{finite transition kernel} from $\Omega_1$ to $\Omega_2$ if:
\begin{enumerate}[label=(\roman*)]
\item for every $A_2\in\A_2$, the function $\omega_1\mapsto \vartheta(\omega_1,A_2)$ is $\A_1$-measurable;
\item for every $\omega_1\in\Omega_1$, the set function $A_2\mapsto \vartheta(\omega_1,A_2)$ is a finite measure on $(\Omega_2,\A_2)$.
\end{enumerate}
If the measures in (ii) are probability measures, $\vartheta$ is a \emph{stochastic/Markov kernel}.
\end{definition}

\begin{definition}[Regular conditional distribution (RCD)]
\label{def:rcd}
Let $Y:(\Omega,\A)\to(E,\E)$ be measurable and let $\F\subseteq \A$ be a sub-$\sigma$-algebra.
A Markov kernel $\vartheta_{Y,\F}$ from $(\Omega,\F)$ to $(E,\E)$ is called a \emph{regular conditional distribution of $Y$ given $\F$} if
\begin{equation}
\label{eq:rcd-def}
\int_A \vartheta_{Y,\F}(\omega,B)\,\mathbb{P}(d\omega) \;=\; \mathbb{P}\big(A\cap \{Y\in B\}\big)
\quad\text{for all }A\in\F,\; B\in\E.
\end{equation}
Equivalently, $\vartheta_{Y,\F}(\cdot,B)$ is a version of $\mathbb{P}(Y\in B\mid \F)$ for every $B\in\E$.
\end{definition}



\paragraph{So what does RCD give us?}
A regular conditional distribution upgrades the abstract identity
\[
\mathbb{E}[\mathbf{1}_{\{Y\in B\}}\mid \F]
\quad\text{(defined only up to $\mathbb{P}$-null sets, separately for each $B$)}
\]
into a single random probability measure
\[
\omega \mapsto \vartheta_{Y,\F}(\omega,\cdot)
\]
such that for almost every $\omega$ we can do \emph{all} of the following \emph{pointwise}:
\begin{itemize}
\item evaluate $\mathbb{P}(Y\in B\mid \F)(\omega)$ for \emph{every} measurable $B$,
\item integrate functions:
\[
\mathbb{E}[f(Y)\mid \F](\omega)=\int f(y)\,\vartheta_{Y,\F}(\omega,dy),
\]
\item treat conditioning as a genuine ``posterior distribution'' of $Y$ given the information $\F$.
\end{itemize}
This is the sense in which regular conditional distributions provide a \emph{regularized} (pointwise, measure-valued) form of conditional expectation.

\subsection*{Existence}
Evidently, a natural question is whether our new object exist in our context. We start by showing the existence in $\mathbb{R}$ and later we generalize for larger spaces. The idea is to construct a distribution in $\mathbb{R}$ that satisfies the RCD definition.
\begin{theorem}[Existence of an RCD for real-valued random variables]
\label{thm:rcd-R}
Let $Y:(\Omega,\A)\to(\mathbb{R},\B(\mathbb{R}))$ be real-valued and $\F\subseteq \A$.
Then there exists a regular conditional distribution $\vartheta_{Y,\F}$ of $Y$ given $\F$.
\end{theorem}

\begin{proof}
\Step{1: conditional distribution function on rationals}
For each $r\in\mathbb{Q}$, pick a version
\[
F(r,\cdot) := \mathbb{P}\big(Y\le r\mid \F\big)
\]
which is $\F$-measurable and takes values in $[0,1]$.

\Step{2: enforce monotonicity outside null sets}
If $r\le s$ then $\Ind_{\{Y\le r\}}\le \Ind_{\{Y\le s\}}$, hence by monotonicity of conditional expectation
there exists $N_{r,s}\in\F$ with $\mathbb{P}(N_{r,s})=0$ such that
\[
F(r,\omega)\le F(s,\omega)\quad\text{for all }\omega\notin N_{r,s}.
\]
Let $N_1:=\bigcup_{r\le s,\; r,s\in\mathbb{Q}}N_{r,s}$; then $\mathbb{P}(N_1)=0$ and for $\omega\notin N_1$ the map
$r\mapsto F(r,\omega)$ is nondecreasing on $\mathbb{Q}$.

\Step{3: enforce right limits and boundary conditions}
Using dominated convergence for conditional expectations, we can choose null sets ensuring:
(i) right limits along $r+1/n$ exist for each rational $r$ outside a null set,
(ii) $\inf_{n\in\mathbb{N}}F(-n,\omega)=0$ and $\sup_{n\in\mathbb{N}}F(n,\omega)=1$ outside a null set.
Let $N_2$ be the union of these null sets, and set $N:=N_1\cup N_2$ so that $\mathbb{P}(N)=0$.

\Step{4: extend to all real $z$ by right-continuous hull}
For $\omega\notin N$ and $z\in\mathbb{R}$ define
\[
\widetilde F(z,\omega):=\inf\{F(r,\omega): r\in\mathbb{Q},\ r>z\}.
\]
Then $\widetilde F(\cdot,\omega)$ is nondecreasing and right-continuous, with limits $0$ at $-\infty$ and $1$ at $+\infty$,
hence it is a distribution function. For $\omega\in N$, set $\widetilde F(\cdot,\omega)$ equal to an arbitrary fixed distribution function.

\Step{5: define the kernel and verify measurability}
For each $\omega$, let $\vartheta(\omega,\cdot)$ be the unique probability measure on $(\mathbb{R},\B(\mathbb{R}))$
with distribution function $\widetilde F(\cdot,\omega)$.
For sets of the form $(-\infty,r]$ with $r\in\mathbb{Q}$ we have
\[
\vartheta(\omega,(-\infty,r])=\widetilde F(r,\omega)=F(r,\omega)\quad\text{for }\omega\notin N,
\]
so $\omega\mapsto \vartheta(\omega,(-\infty,r])$ is $\F$-measurable.
Since $\{(-\infty,r]:r\in\mathbb{Q}\}$ generates $\B(\mathbb{R})$ and is stable under finite intersections,
a standard monotone class/$\pi$-$\lambda$ argument yields $\F$-measurability of $\omega\mapsto \vartheta(\omega,B)$ for all Borel $B$.

\Step{6: verify the defining property \eqref{eq:rcd-def}}
For $A\in\F$ and $B=(-\infty,r]$ with $r\in\mathbb{Q}$,
\[
\int_A \vartheta(\omega,B)\,\mathbb{P}(d\omega)=\int_A F(r,\omega)\,\mathbb{P}(d\omega)
= \mathbb{P}\big(A\cap\{Y\le r\}\big)
\]
by the defining property of conditional expectation.
Both sides are finite measures in $B$ and agree on the generator $\{(-\infty,r]: r\in\mathbb{Q}\}$, hence they agree on all $B\in\B(\mathbb{R})$.
Thus $\vartheta$ is a regular conditional distribution of $Y$ given $\F$.
\end{proof}

\section{Borel spaces, Polish spaces, and Lusin spaces}

Regular conditional distributions are \emph{pointwise} objects: for almost every outcome
we want a genuine probability measure that depends on the conditioning information in a \emph{measurable} way.
This measurability requirement is subtle. In ``well-behaved'' measurable spaces one can select such conditional laws;
in pathological spaces this selection may fail. The classes below are precisely the standard ``well-behaved'' settings
in which conditional probabilities behave as one expects from $\mathbb{R}^n$.

\subsection*{Borel spaces}
A Borel space is, by definition, a measurable space that can be relabeled into a Borel subset of $\mathbb{R}$ equipped with its Borel $\sigma$-algebra.
This might look restrictive, but it covers essentially all state spaces used in probability theory:
\[
\begin{aligned}
&\mathbb{R}^n,\quad \mathbb{R}^\infty \text{ (countable product),}\;\text{separable Banach/Hilbert spaces with their Borel $\sigma$-algebras,}\\
&\text{countable sets with the power set $\sigma$-algebra, and any separable complete metric space.}
\end{aligned}
\]

Intuitively, ``standard Borel'' means: there is a \emph{countable amount of information} that generates all measurable sets
(a countable generator), and this makes it possible to construct conditional probabilities by first defining them on the generator
and then extending them to all measurable sets. Essentially, these are abstract spaces that behave like Borel subsets of $\mathbb{R}$.

\begin{definition}[Isomorphism of measurable spaces and measure spaces]
\label{def:isom}
Measurable spaces $(E,\E)$ and $(\widetilde E,\widetilde{\E})$ are \emph{isomorphic}
if there exists a bijection $\xi:E\to \widetilde E$ such that $\xi$ is $\E$-$\widetilde{\E}$ measurable and
$\xi^{-1}$ is $\widetilde{\E}$-$\E$ measurable.
If moreover $\mu$ on $(E,\E)$ and $\widetilde\mu$ on $(\widetilde E,\widetilde{\E})$ satisfy
$\widetilde\mu=\mu\circ \xi^{-1}$, then $(E,\E,\mu)$ and $(\widetilde E,\widetilde{\E},\widetilde\mu)$ are
\emph{isomorphic as measure spaces}.
\end{definition}

\begin{definition}[Borel space / standard Borel space]
\label{def:borelspace}
A measurable space $(E,\E)$ is called a \emph{Borel space} (also: \emph{standard Borel})
if it is isomorphic to $(B,\B(B))$ for some Borel set $B\in \B(\mathbb{R})$.
\end{definition}

\subsection*{Polish space}
A Polish space is a separable complete metric space.
This is a topological condition, not a purely measurable one, but it is the most common source of standard Borel spaces:
if $E$ is Polish, then its Borel $\sigma$-algebra $\B(E)$ is standard Borel (see \ref{thm:polish-to-borel}).
Concrete parallels are:
\[
\mathbb{R}^n,\quad \text{any separable Hilbert space }(\ell^2, L^2[0,1]),\quad C([0,1]),\quad \text{separable Banach spaces}.
\]
The key intuition is that separability provides a countable dense set (hence countable ``coordinates'' to approximate open sets),
and completeness ensures limits of Cauchy sequences stay inside the space. Together, these properties yield a robust notion of
approximation and enable powerful selection and regularity theorems in probability (existence of regular conditional distributions,
measurable selections, tightness arguments, etc.).

\begin{definition}[Polish space]
A topological space is \emph{Polish} if it is separable and its topology is induced by a complete metric.
\end{definition}

\begin{theorem}[Polish $\Rightarrow$ Borel space]
\label{thm:polish-to-borel}
If $E$ is Polish and $\E=\B(E)$ is its Borel $\sigma$-algebra, then $(E,\E)$ is a Borel space.
\end{theorem}

\subsection*{Lusin spaces}

Lusin spaces are often presented as spaces that are, in a precise sense, not necessarily Polish themselves,
but can be realized as a measurable image of a Polish space without losing the ``good'' measurability properties. They are ``almost Polish'' from the measurable viewpoint.
A typical formulation is: a Lusin measurable space is isomorphic (as a measurable space) to a Borel subset of a Polish space.
Thus, from the standpoint of measure theory, Lusin spaces retain the same good behavior as standard Borel spaces.
A helpful parallel is manifold theory: a manifold may not be globally Euclidean, but locally it looks like $\mathbb{R}^n$;
similarly, a Lusin space may not come with a nice global metric structure, but it still carries the same measurable structure
as a Borel subset of a Polish space. 

\begin{definition}[Lusin measurable space]
A measurable space $(E,\E)$ is called a \emph{Lusin measurable space} if there exists a Polish space $P$
and a Borel set $B\in\B(P)$ such that $(E,\E)$ is isomorphic (as a measurable space) to $(B,\B(B))$.
In other words, there exists a bijection $\varphi:E\to B$ with $\varphi$ and $\varphi^{-1}$ measurable.
\end{definition}

\subsection{Existence of RCDs on Borel spaces}

\begin{theorem}[Existence of RCDs for Borel-valued random variables]
\label{thm:rcd-borel}
Let $\F\subseteq\A$ and let $Y:(\Omega,\A)\to(E,\E)$ take values in a Borel space $(E,\E)$.
Then there exists a regular conditional distribution $\vartheta_{Y,\F}$ of $Y$ given $\F$.
\end{theorem}

\begin{proof}
Since $(E,\E)$ is a Borel space, there exist a Borel set $B\subseteq\mathbb{R}$ and a measurable isomorphism
$\xi:(E,\E)\to(B,\B(B))$.
Define the real-valued random variable $\widetilde Y := \xi\circ Y$ (viewed as a map into $\mathbb{R}$ by extending outside $B$ arbitrarily).
By Theorem~\ref{thm:rcd-R}, there exists a regular conditional distribution $\vartheta_{\widetilde Y,\F}$ of $\widetilde Y$ given $\F$.
Now define, for $A\in\E$,
\[
\vartheta_{Y,\F}(\omega,A):=\vartheta_{\widetilde Y,\F}\big(\omega,\xi(A)\big).
\]
Measurability in $\omega$ and the measure property in $A$ follow from those of $\vartheta_{\widetilde Y,\F}$ and the isomorphism.
Finally, for $C\in\F$,
\[
\int_C \vartheta_{Y,\F}(\omega,A)\,\mathbb{P}(d\omega)
=\int_C \vartheta_{\widetilde Y,\F}(\omega,\xi(A))\,\mathbb{P}(d\omega)
=\mathbb{P}\big(C\cap\{\widetilde Y\in \xi(A)\}\big)
=\mathbb{P}\big(C\cap\{Y\in A\}\big),
\]
so $\vartheta_{Y,\F}$ is an RCD.
\end{proof}

\subsection{Regularized conditional expectation: kernel representation}

The next statement explains why a regular conditional distribution is a \emph{regularization} of conditional expectation:
it turns $\mathbb{E}[f(X)\mid \F]$ into an integral against a (random) probability measure.

\begin{theorem}[Integral representation of conditional expectation]
\label{thm:reg-cond-exp}
Let $X:(\Omega,\A)\to(E,\E)$ be Borel-valued and let $\F\subseteq\A$.
Let $\vartheta_{X,\F}$ be an RCD of $X$ given $\F$.
If $f:E\to\mathbb{R}$ is measurable with $\mathbb{E}[|f(X)|]<\infty$, then
\begin{equation}
\label{eq:kernel-repr}
\mathbb{E}[f(X)\mid \F](\omega)=\int_E f(x)\,\vartheta_{X,\F}(\omega,dx)
\qquad\text{for }\mathbb{P}\text{-a.e. }\omega.
\end{equation}
\end{theorem}

\begin{proof}
It suffices to consider $f\ge 0$ (apply to $f^+$ and $f^-$ otherwise).
Approximate $f$ increasingly by simple functions:
there exist measurable sets $A_i\in\E$ and $\alpha_i\ge 0$ such that
\[
g_n:=\sum_{i=1}^n \alpha_i \Ind_{A_i}\uparrow f.
\]
Fix $B\in\F$. Using Definition~\ref{def:rcd} with $Y=X$,
\begin{align*}
\mathbb{E}[g_n(X)\Ind_B]
&=\sum_{i=1}^n \alpha_i\,\mathbb{P}(B\cap\{X\in A_i\})
=\sum_{i=1}^n \alpha_i \int_B \vartheta_{X,\F}(\omega,A_i)\,\mathbb{P}(d\omega) \\
&=\int_B \left(\int_E g_n(x)\,\vartheta_{X,\F}(\omega,dx)\right)\mathbb{P}(d\omega).    
\end{align*}

By monotone convergence, for $\mathbb{P}$-a.e.\ $\omega$,
$\int g_n(x)\,\vartheta_{X,\F}(\omega,dx)\uparrow \int f(x)\,\vartheta_{X,\F}(\omega,dx)$,
and also $\mathbb{E}[g_n(X)\Ind_B]\uparrow \mathbb{E}[f(X)\Ind_B]$.
Thus, for all $B\in\F$,
\[
\mathbb{E}[f(X)\Ind_B]
=\int_B \left(\int_E f(x)\,\vartheta_{X,\F}(\omega,dx)\right)\mathbb{P}(d\omega).
\]
The right-hand side is $\F$-measurable in $\omega$, so it has the defining property of $\mathbb{E}[f(X)\mid \F]$.
\end{proof}

\section{Disintegration lemma}
A common situation one faces is to \emph{simulate} or \emph{integrate} with respect to a joint model $(X,Y)$ whose law $\mu$ on $E\times F$ is complicated, while the marginal of $X$ and the conditional law of $Y$ given $X=x$ are tractable.  The disintegration lemma justifies rigorously the intuitive two-step procedure: \emph{first sample $x$ from the marginal law $\mu_E$ of $X$, and then sample $y$ from the conditional law $\kappa(x,\cdot)$ given this $x$.}  In this way, probabilities and expectations under the joint law can be computed by iterating: for instance, for an integrable payoff $\varphi:E\times F\to\mathbb{R}$ one can evaluate
\[
\int_{E\times F}\varphi(x,y)\,\mu(dx,dy)
=\int_E\left(\int_F \varphi(x,y)\,\kappa(x,dy)\right)\mu_E(dx),
\]
which is exactly the mathematical statement behind many practical algorithms (e.g.\ hierarchical simulation, Monte--Carlo estimators, Bayesian posterior prediction).  The key point is that this remains valid even when there is \emph{no density} and no ``formula'' for $\kappa(x,dy)$: disintegration guarantees that such a measurable conditional kernel exists and that the two-step computation truly reproduces the original joint law.

\paragraph{Idea.}
A probability measure $\mu$ on a product space $E\times F$ is the mathematical object that encodes the
\emph{joint law} of a pair of random variables $(X,Y)$ taking values in $E$ and $F$.
From basic probability one expects that a joint law should decompose into
\[
\text{(law of }X)\times \text{(conditional law of }Y\text{ given }X=x).
\]
This is exactly what the disintegration lemma formalizes.

More concretely, let $\mu_E$ be the first marginal of $\mu$, i.e.\ the distribution of $X$.
Then disintegration asserts that there exists a measurable family of probability measures
\[
x \longmapsto \kappa(x,\cdot)\in \mathcal{P}(F),
\]
such that for each measurable set $B\subseteq F$, the map $x\mapsto \kappa(x,B)$ is measurable and
\[
\mu(A\times B)=\int_A \kappa(x,B)\,\mu_E(dx)\qquad (A\in\E,\;B\in\F).
\]
Heuristically, you should read $\kappa(x,\cdot)$ as the \emph{conditional distribution} of $Y$ given $X=x$.
Then the identity above says: to compute the probability that $(X,Y)$ falls in $A\times B$,
first choose $x$ according to the law of $X$, and then (given this $x$) choose $y$ according to the conditional law of $Y$ given $X=x$.

\paragraph{Why this is not automatic from conditional expectation?}
Conditional expectation guarantees that for each fixed set $B\in\F$ the random variable
$\mathbb{P}(Y\in B\mid \sigma(X))$ exists (as an a.s.\ defined object).
Disintegration goes further: it produces \emph{one single} object $\kappa(x,\cdot)$ such that, for almost every $x$,
$B\mapsto \kappa(x,B)$ is a genuine probability measure and the dependence on $x$ is measurable.
In other words, the lemma upgrades ``separate conditional probabilities'' into a \emph{measurable kernel}
(which can then be used to integrate arbitrary functions of $Y$ given $X=x$).






\begin{theorem}[Disintegration lemma for probability measures on products]
\label{thm:disintegration}
Let $(E,\E)$ and $(F,\F)$ be Borel spaces and let $\mu$ be a probability measure on $(E\times F,\E\otimes \F)$.
Let $\mu_E$ be the first marginal: $\mu_E(A):=\mu(A\times F)$.
Then there exists a Markov kernel $\kappa:E\times \F\to[0,1]$ such that:
\begin{enumerate}[label=(\roman*)]
\item for each $B\in\F$, the map $x\mapsto \kappa(x,B)$ is $\E$-measurable;
\item for each $x\in E$, $B\mapsto \kappa(x,B)$ is a probability measure on $(F,\F)$;
\item (\emph{disintegration identity}) for all $A\in\E$ and $B\in\F$,
\begin{equation}
\label{eq:disintegration}
\mu(A\times B)=\int_A \kappa(x,B)\,\mu_E(dx).
\end{equation}
\end{enumerate}
Moreover, $\kappa$ is unique $\mu_E$-a.e.\ in the sense that if $\kappa'$ also satisfies \eqref{eq:disintegration},
then $\kappa(x,\cdot)=\kappa'(x,\cdot)$ for $\mu_E$-a.e.\ $x$.
\end{theorem}

\begin{proof}
Consider the probability space $(\Omega,\A,\mathbb{P}):=(E\times F,\E\otimes\F,\mu)$.
Let $X(\omega)=x$ and $Y(\omega)=y$ be the coordinate projections on $E$ and $F$, respectively.
Then $X$ takes values in the Borel space $(E,\E)$ and $Y$ in $(F,\F)$.

By Theorem~\ref{thm:rcd-borel}, there exists a regular conditional distribution
$\vartheta_{Y,\sigma(X)}$ of $Y$ given $\sigma(X)$.
Since $\sigma(X)$ consists of sets of the form $X^{-1}(A)=A\times F$ with $A\in\E$,
the kernel $\vartheta_{Y,\sigma(X)}(\omega,\cdot)$ is $\sigma(X)$-measurable in $\omega$.

\Step{1: factor the $\sigma(X)$-measurable kernel through $X$}
Fix $B\in\F$. The map $\omega\mapsto \vartheta_{Y,\sigma(X)}(\omega,B)$ is $\sigma(X)$-measurable, hence by the factorization Lemma
there exists a measurable map $k_B:E\to[0,1]$ such that
\[
\vartheta_{Y,\sigma(X)}(\omega,B)=k_B(X(\omega))\qquad \mu\text{-a.s.}
\]
Define $\kappa(x,B):=k_B(x)$. One checks that $x\mapsto \kappa(x,B)$ is measurable for each $B$.

\Step{2: $B\mapsto \kappa(x,B)$ is a probability measure}
For a.e.\ $\omega$, $B\mapsto \vartheta_{Y,\sigma(X)}(\omega,B)$ is a probability measure.
Since $\vartheta_{Y,\sigma(X)}(\omega,\cdot)=\kappa(X(\omega),\cdot)$ a.s., it follows that for $\mu_E$-a.e.\ $x$,
$B\mapsto \kappa(x,B)$ is a probability measure. Redefine $\kappa(x,\cdot)$ arbitrarily on a $\mu_E$-null set of $x$ to make it a probability measure for all $x$.

\Step{3: verify the disintegration identity}
Take $A\in\E$ and $B\in\F$. Using the RCD property \eqref{eq:rcd-def} with $\F=\sigma(X)$ and $A\times F\in\sigma(X)$,
\[
\mu(A\times B)
= \int_{A\times F} \vartheta_{Y,\sigma(X)}(\omega,B)\,\mu(d\omega)
= \int_{A\times F} \kappa(X(\omega),B)\,\mu(d\omega).
\]
But under $\mu$, $X(\omega)=x$ and integrating over $A\times F$ amounts to integrating $x\in A$ against the marginal $\mu_E$:
\[
\int_{A\times F} \kappa(X(\omega),B)\,\mu(d\omega)=\int_A \kappa(x,B)\,\mu_E(dx),
\]
which is \eqref{eq:disintegration}.

\LabelProof{Uniqueness}
If $\kappa'$ also satisfies \eqref{eq:disintegration}, then for each fixed $B\in\F$,
the functions $x\mapsto \kappa(x,B)$ and $x\mapsto \kappa'(x,B)$ have the same integrals over all $A\in\E$ w.r.t.\ $\mu_E$.
Hence they are equal $\mu_E$-a.e. for each $B$ in a countable generator of $\F$.
A monotone class argument extends equality to all $B\in\F$, yielding $\kappa(x,\cdot)=\kappa'(x,\cdot)$ for $\mu_E$-a.e.\ $x$.
\end{proof}

\begin{remark}[Finite measures and $\sigma$-finite extensions]
If $\mu$ is a finite measure, apply Theorem~\ref{thm:disintegration} to $\mu/\mu(E\times F)$ and rescale.
Extensions to $\sigma$-finite measures require additional hypotheses; for standard Borel spaces one typically assumes a suitable $\sigma$-finite structure
to ensure existence of a disintegration kernel.
\end{remark}

\begin{remark}{The geometric interpretation}
You can think of the product space $E\times F$ as being ``foliated'' by the vertical fibers
\[
\{x\}\times F,\qquad x\in E.
\]
The joint measure $\mu$ assigns mass to sets in the product.
Disintegration says that one can ``slice'' $\mu$ along these fibers:
for (almost every) $x$ there is a conditional measure $\kappa(x,\cdot)$ on the fiber $\{x\}\times F$,
and the original mass is recovered by averaging these slices with respect to the marginal $\mu_E$.
This is the measure-theoretic analogue of Fubini's theorem:
Fubini integrates a function by iterated integrals, while disintegration decomposes the \emph{measure itself}
into an iterated structure (first in $x$, then in $y$ conditioned on $x$).
\end{remark}

\begin{remark}{An analogy to densities.}
If $\mu$ admits a joint density $p_{X,Y}(x,y)$ with respect to Lebesgue measure, then one writes
\[
p_{X,Y}(x,y)=p_X(x)\,p_{Y\mid X}(y\mid x),
\]
where $p_{Y\mid X}(y\mid x)=\frac{p_{X,Y}(x,y)}{p_X(x)}$ whenever $p_X(x)>0$.
The disintegration lemma is the \emph{measure-level} version of this factorization,
and it remains valid even when densities do \emph{not} exist (discrete/continuous mixtures, singular laws, etc.).
Thus, it is the correct general replacement for the ``ratio formula'' for conditional densities.
\end{remark}